%\chapter{Reference guide}

\section{Toroidal momentum transport equation}
In this section the newly implemented transport equation to advance toroidal 
momentum is derived in details and the options to call it in the model file are
presented. 

The motivation that has pushed to add this important piece of physics stems from
the fact that application of toroidal momentum transport theory with respect to 
experimental studies has rapidly expanded and it is now a very important part of
the fusion research. In the previous Astra versions, it was possible to mock--up
momentum transport via the F{\it i} equations. However that has strong 
limitations and the physical representation of momentum transport through 
density equations is not accurate in most cases. That is why this equation has
been added to the code.

The derivation presented here follows standard literature, from old classic derivations as in Hinton--Hazeltine review on tokamak plasma transport theory, to more recent as in V. Rozhansky {\it et al}. or in J. D. Callen {\it et al}.
%
\subsection{Derivation of the equation implemented in Astra--7.0}
The basic assumption behind the derivation of the toroidal momentum transport equation is macroscopic plasma quasi--neutrality:
\begin{eqnarray}
\nabla \cdot \mathbf{J} = 0
\label{qneutraldivjeq0}
\end{eqnarray}
Its flux--surface ($\displaystyle \langle f \rangle \doteq \frac{\partial}{\partial{V}}\int{f dV}$) average reads 
\begin{eqnarray}
\langle J_\Psi R B_p \rangle=0
\label{qneutraldivjeq0fsa}
\end{eqnarray}
with $J_r$ the 
net radial current of the bulk plasma, $R$ the local major radius, and $B_p$ the
poloidal magnetic field.

The radial current density is given by non--ambipolar particle transport. In the case of a multi--species plasma this writes down as:
\begin{eqnarray}
J_\Psi =e\sum_s Z_s \Gamma_s 
\label{jpsimultispecpartsum}
\end{eqnarray}
with $e Z_s$ the species charge, and $\Gamma_s$ the cross--flux surface particle flux of species $s$. 
%
\subsection{Evaluation of the radial current and consequences on momentum transport}
Each species $s$ has its own momentum transport equation, that is the complete three--dimensional force balance:
\begin{eqnarray}
\frac{\partial(m_s n_s\Uvec_s)}{\partial{t}}+\nabla\cdot\left(\Pivec_{c,s}
+\Pivec_{n,s}+\Pivec_{a,s}\right)=e Z_s n_s \left(\Evec+\Uvec_s\times\Bvec\right) +\Fvec_s+\Svec_s
\label{momtranspeq3d}
\end{eqnarray}
The different quantities and terms are respectively: $m_s$ the species mass, $n_s$ the species density, $\Pivec_{x,s}$ 
the species anisotropic stress tensor (c=classical, n=neoclassical, a=anomalous). Note that the anomalous part 
supposedly contains also the turbulent part of the Maxwell stress tensor. $\Evec,\Bvec$ are the electric and magnetic 
fields, $\Fvec_s$ the friction force, and $\Svec_s$ the externally applied torque (NBI for example).

We now multiply scalarly \eiref{momtranspeq3d} by $R\ephi$ and sum over all species:
\begin{eqnarray}
 \frac{\partial M_\phi}{\partial{t}}+\sum_s R\ephi\cdot
\nabla\cdot\left(\Pivec_{c,s}
+\Pivec_{n,s}+\Pivec_{a,s}\right)=R\ephi\cdot\left(
\Jvec\times\Bvec\right) + R \sum_s S_{\phi,s}
\label{momtranspeq3drephi}
\end{eqnarray}
where $M_{\phi}=\sum_s m_s n_s R U_{\phi,s}$ is the plasma total toroidal angular momentum density.
Notice that, for an axisymmetric system, $R\ephi\cdot(\nabla\cdot) \rightarrow \nabla\cdot (R \ephi\cdot)$, leading to:
\begin{eqnarray}
 \frac{\partial M_\phi}{\partial{t}}+
\sum_s \nabla\cdot\left(R\ephi\cdot\Pivec_s\right)=R\ephi\cdot\left(
\Jvec\times\Bvec\right)  + R \sum_s S_{\phi,s}
\label{momtranspeq3drephi2}
\end{eqnarray}
where $\Pivec_s=\Pivec_{c,s}+\Pivec_{n,s}+\Pivec_{a,s}$
On the other hand: $R\ephi\cdot\left(
\Jvec\times\Bvec\right)=R J_\psi \bpol$. 

Before taking the flux--surface average, we rewrite \eiref{momtranspeq3drephi2} in 'conservative--like' form:
\begin{eqnarray}
\int_V\left[\frac{\partial M_\phi}{\partial{t}} +
\sum_s \nabla\cdot\left(R\ephi\cdot\Pivec_s\right)-R\ephi\cdot\left(
\Jvec\times\Bvec\right)  - R \sum_s S_{\phi,s}\right] dV=0
\label{momtranspeq3drephi2cons}
\end{eqnarray}
The volume integral spans inside a closed flux surface identified by a constant label $\Phi$, which we choose as the toroidal flux enclosed by said surface. In that case the integral can be swapped with the time differential through the operation:
\begin{eqnarray}
 \int_V \frac{\partial M_\phi}{\partial{t}} dV=\left[\frac{\partial }{\partial t}  \int_V M_\phi dV \right]_{\Phi} -\oint_{\Phi} M_\phi \vvec_{\Phi}\cdot \nabla{V} \frac{dS}{|\nabla{V}|}
\label{momtranspeq3drephi2consswap}
\end{eqnarray}
By doing this, and differentiating then by $\Phi$ at constant time, we get the equation ($V'=\partial V/\partial \Phi$):
\begin{eqnarray}
\frac{1}{V'}  \frac{\partial }{\partial{t}}_{\Phi}\left(V' \left\langle M_\phi \right\rangle\right) +\frac{\partial}{\partial V}\left\langle \left(R \sum_s \Pi_{\phi\Psi,s}-M_\phi\vvec_{\Phi}\cdot\frac{\nabla{V}}{|\nabla{V}|}\right)
|\nabla{V}| \right\rangle\nonumber \\
-\left\langle R J_\Psi\bpol\right\rangle  - 
\left\langle R \sum_s S_{\phi,s}\right\rangle=0
\label{momtranspeq3drephi2cons}
\end{eqnarray}
Therefore, imposing condition (\ref{qneutraldivjeq0fsa}), we end up with the relevant equation:
\begin{eqnarray}
\frac{1}{V'}  \frac{\partial }{\partial{t}}_{\Phi}\left(V' \left\langle M_\phi \right\rangle\right) + \frac{\partial}{\partial V}\left\langle \left(R \sum_s \Pi_{\phi\Psi,s}-M_\phi\vvec_{\Phi}\cdot\frac{\nabla{V}}{|\nabla{V}|}\right)
|\nabla{V}| \right\rangle=\nonumber \\ 
=\left\langle R \sum_s S_{\phi,s}\right\rangle
\label{momtranspeq3drephiFSA}
\end{eqnarray}
The equation can also be written in terms of the coordinate $\rho=\sqrt{\Phi/(\pi B_0)}$ that is used in ASTRA ($V'=\partial V/\partial \rho$), and $B_0=\Bvacr$ is the reference {\bf BTOR} field:
\begin{eqnarray}
\frac{1}{V'} \frac{\partial }{\partial{t}}_{\rho}\left(V' \left\langle M_\phi \right\rangle\right) +\frac{1}{V'}  \frac{\partial}{\partial \rho}\left[ V' \left\langle R \sum_s \Pi_{\phi\Psi,s}\vert \nabla{\rho}\vert\right\rangle
-V'\rho\frac{\dot{B}_0}{2 B_0}\langle  M_\phi \rangle \right]
=\nonumber \\ 
=\left\langle R \sum_s S_{\phi,s}\right\rangle
\label{momtranspeq3drephiFSArho}
\end{eqnarray}
or
\begin{eqnarray}
\frac{1}{V'} \frac{\partial }{\partial{t}}_{\rho}\left(V' \left\langle M_\phi \right\rangle\right) +\frac{1}{V'}  \frac{\partial}{\partial \rho}\left[ V' R_{\phi}
-V'\rho\frac{\dot{B}_0}{2 B_0}\langle  M_\phi \rangle \right]
=\nonumber \\ 
=\left\langle R \sum_s S_{\phi,s}\right\rangle
\label{momtranspeq3drephiFSArho}
\end{eqnarray}

Now, we have to express $M= \left\langle M_\phi \right\rangle$ in some practical way. Since $\displaystyle  M_\phi = R \sum_s m_s n_s U_{\phi,s}$, we write $U_{\phi,s}=R\Omega_s$, where $\Omega_s$ is the angular rotation frequency.

We have to find a way to express the multispecies velocity, and to express the rotation as a flux function.

Let us start from (for a generic species): $\Uvec=U_\parallel \bvec+\Uperpvec$, and $\Uvec=U_\phi \ephi + U_p \eP$, where
$\displaystyle \Uperpvec=\frac{\left[\Evec-\nabla{P}/(Zen)\right]\times\Bvec}{B^2}$.
Note that $H=(E_r-\partial_r P /(Zen))/(R \bpol)=H_{\rm E}+H_{\rm P}$ (i.e. split in electric field and pressure contributions) is a flux function. Also, $U_p/\bpol=K(\Psi)$ is a flux function.

So, we get:
\begin{eqnarray}
U_\parallel = K B + \frac{I}{B} (H_{\rm E}+H_{\rm P})
\label{uparperp1}
\end{eqnarray}
and
\begin{eqnarray}
U_\phi = K B_\phi + R (H_{\rm E}+H_{\rm P})
\label{uparperp1}
\end{eqnarray}
or
\begin{eqnarray}
 I H_{\rm E}=U_\parallel B-I  H_{\rm P}-K B^2
\label{uparperp2}
\end{eqnarray}
Note that $H = H(\Psi)$, i.e. is a flux function, since $\displaystyle H(\Psi)=-\left[\frac{\partial\Phi_0}{\partial \Psi}+\frac{1}{Zen}\frac{\partial{P}}{\partial{\Psi}}\right]$. 

Let us define 
\begin{eqnarray}
\upara=\left\langle U_\parallel B \right\rangle/B_0 = K \frac{\left\langle B^2 \right\rangle}{B_0}+\frac{I}{B_0}(H_{\rm E}+H_{\rm P})
\label{uparfsadef}
\end{eqnarray}
and
\begin{eqnarray}
 I H_{\rm E}=B_0 \upara-I  H_{\rm P}-K \left\langle B^2 \right\rangle
\label{uparfsadef2}
\end{eqnarray}

so that:
\begin{eqnarray}
U_\phi = R \frac{B_0}{I}\upara + K I \left(\frac{1}{R}-R\frac{\left\langle B^2 \right\rangle}{I^2}\right)
\label{omegafromupar}
\end{eqnarray} 
Then:
\begin{eqnarray}
\Omega =\frac{B_0}{I}\upara + K I \left(\frac{1}{R^2}-\frac{\left\langle B^2 \right\rangle}{I^2}\right)
\label{uparperp4}
\end{eqnarray}
This last equation allows to relate the code result for $\upara$ to the experimental measurement of $\Omega$ at some location. 

So, we have to compute $M= \left\langle M_\phi \right\rangle$, which, neglecting poloidal variations of the density (this will be upgraded in a future version), can be written as:
\begin{eqnarray}
\left\langle R^2 \Omega \right\rangle = \frac{\left\langle R^2 \right\rangle B_0}{I}\upara+KI\left(1- \frac{\langle R^2  \rangle \langle B^2 \rangle }{I^2} \right)
\label{r2omupara}
\end{eqnarray}

In addition, for each impurity species $s\neq D$, the neoclassical formulas provide $\uparaDZ=\upara^{\rm D}-\upara^{\rm s}$.
So that we can write:
\begin{eqnarray}
\sum_s m_s n_s \upara^{s}=\rho_{\rm m} \upara-\sum_s m_s n_s \uparaDZ
\label{mphineoc}
\end{eqnarray}
with the total mass density $\rho_{\rm m}=\sum_s m_s n_s $. We also define $\delta_{\rm neo}=\sum_s m_s n_s \uparaDZ$, and $\sigma_{\rm neo}=\sum_s m_s n_s K_s$.

We can then rewrite the average toroidal mass flow as:
\begin{eqnarray}
\left\langle M_\phi \right\rangle=\Upsilon_0 \upara + \Upsilon_1 + \Upsilon_2
\label{mphiavupar}
\end{eqnarray}
with 
\begin{eqnarray}
\Upsilon_0&=&\rho_{\rm m}\frac{\left\langle R^2 \right\rangle\Bvacr}{I}
\nonumber \\
\Upsilon_1&=&-\frac{\left\langle R^2 \right\rangle\Bvacr}{I}\delta_{\rm neo}
\nonumber \\
\Upsilon_2&=&I\left(1- \frac{\langle R^2  \rangle \langle B^2 \rangle }{F} \right)\sigma_{\rm neo} 
\label{ABdefs}
\end{eqnarray}
Note that in the limit of zero poloidal velocity and infinite friction, we can simplify $\Upsilon_1=\Upsilon_2=0$.

The choice of $\upara$ as the variable to be solved for is that the parallel direction is natural for plasma transport, as such theory--based transport model provide transport of parallel velocity.
Infact, the term $\displaystyle \left\langle R \sum_s \Pi_{\phi\Psi,s}
|\nabla{\rho}| \right\rangle$ can be represented by considering that:
\begin{eqnarray}
 \Pi_{\phi\Psi,s} \rightarrow  \Pi_{\phi\Psi,s}+\frac{\bphi}{B}\Pi_{\parallel\Psi,s}-\frac{\bpol}{B}\Pi_{y\Psi,s}+m_s \Uphis \Gamma_s
\label{stresstensorparperp}
\end{eqnarray} 
where the parallel and binormal stress--tensors represent viscous terms at a fixed $\Phi$ surface. Also here $\Gamma_s$ is the particle flux.
Also:
\begin{eqnarray}
R_\phi \rightarrow  R_\phi+R_\parallel-R_y+m_s \Uphis \Gamma_s
\label{stresstensorparperp}
\end{eqnarray} 



The parallel stress tensor is in general given by (the species $s$ is intended, and also at the moment the residual stress terms are not included):
\begin{eqnarray}
\left\langle R \frac{\bphi}{B} \Pi_{\parallel\Psi}
|\nabla{\rho}| \right\rangle=-m n R_0^2\left\langle |\nabla{\rho}|^2 \right\rangle \left[\chi \left(\frac{\partial H_{\rm e}}{\partial\rho}+\frac{1}{I}\frac{\partial (I H_{\rm P})}{\partial\rho}\right)+\chi_\theta \frac{b_2}{I}\frac{\partial K}{\partial\rho}-V H_{\rm e}-V_\theta \frac{b_2}{I} K -V_d H_{\rm P}\right]
\label{parallelstresstensor}
\end{eqnarray} 
where $\chi,\chi_\theta$ are the turbulent viscoties of the toroidal and poloidal flow respectively, $V,V_\theta,V_d$ are the pinch of electric, poloidal, and diamagnetic flows. $b_2=\left\langle B^2 \right\rangle$.
The expression can be rewritten upon substituting using formulas (\ref{uparfsadef2}), and summing over species:
\begin{eqnarray}
\left\langle R \frac{\bphi}{B} \Pi_{\parallel\Psi}
|\nabla{\rho}| \right\rangle= \frac{R_0^2 B_0}{I} \left\langle |\nabla{\rho}|^2 \right\rangle \left[-a_1 \frac{\partial \upara}{\partial\rho}+a_2\upara+a_3\delta_{\rm neo}'+a_4\delta_{\rm neo} +a_5 \sigma_{\rm neo}'+a_6\sigma_{\rm neo}+a_7 \delta_{\rm d,neo}\right]
\label{parallelstresstensor2}
\end{eqnarray} 
with the following definitions:
\begin{eqnarray}
a_1 &=& \rho_{\rm m} \chi
\nonumber \\
a_2 &=& \rho_{\rm m} \left(V+\chi \frac{\partial \log{I}}{\partial \rho}  \right)
\nonumber \\
a_3 &=& \chi
\nonumber \\
a_4 &=& -\left(V+\chi \frac{\partial \log{I}}{\partial \rho}  \right)
\nonumber \\
a_5 &=& \frac{b_2}{B_0}\left(\chi-\chi_\theta \right)
\nonumber \\
a_6 &=& \frac{b_2}{B_0}\left(V_\theta-V-\chi \frac{\partial \log{I}}{\partial \rho}+\chi\frac{\partial \log{b_2}}{\partial \rho} \right)
\nonumber \\
a_7 &=& \left(V_d-V-\chi\frac{\partial \log{I}}{\partial \rho}  \right)
\nonumber \\
\delta_{\rm neo}' &=& \sum_s m_s n_s \frac{\partial \uparaDZ}{\partial\rho}
\nonumber \\
\sigma_{\rm neo}' &=& \sum_s m_s n_s \frac{\partial K_s}{\partial\rho} 
\nonumber \\
\delta_{\rm d,neo} &=& \sum_s m_s n_s \frac{I H_{\rm P}}{B_0}  
\label{stresstensorparperp2quantities}
\end{eqnarray} 
Of course all of the above assumes that the transport coefficients are the same for all ion species (this will be improved in a future upgrade).

The binormal stress tensor can be kept as a whole since in general cannot be represented in terms of diffusion and convection.
%
\subsection{Internal variables}
The parallel velocity to be solved for would be:
\begin{eqnarray}
{\rm UPAR} = \upara
\label{upardefastra}
\end{eqnarray}
The main ions diffusivity, pinch, and residual stresses (assumed to be the same for all ion species at the moment) are
given in the table below:
\begin{eqnarray}
{\rm XUPAR} &=& \chi \nonumber \\
{\rm XUPAP} &=& \chi_\theta \nonumber \\
{\rm CNPAR} &=& V \nonumber \\
{\rm CNPAD} &=& V_d \nonumber \\
{\rm CNPAP} &=& V_\theta \nonumber \\
{\rm RUPAR} &=& R_{\parallel} \nonumber \\
{\rm RUPYR} &=& R_y  \nonumber \\
{\rm RUPFR} &=& R_{\phi} \nonumber \\
{\rm TTRQ} &=& \left\langle R \sum_s S_{\phi,s}\right\rangle
\label{turbstuffdefastra}
\end{eqnarray}
\\
{\bf Table 2.1} Variables (arrays) related to the momentum transport equation\\[2ex]
\begin{tabular}{|l|l|l|l|}					\hline
\bf Name	&\bf Units & Definition & Description\\ \hline
\tt MRHO 	& kg/m$^3/($10^{19}$ $m_{\rm proton}$) $	& $\rho_{\rm m}$ & Total mass density $\rho_{\rm m}=\sum_s m_s n_s $\\
\tt DLNEO 	& kg/(m$^2$ s)/($10^{19}$ $m_{\rm proton}$) & $\delta_{\rm neo}$ & $\delta_{\rm neo}=\sum_s m_s n_s \uparaDZ$ \\
\tt DLNEOD 	& kg/(m$^3$ s)/($10^{19}$ $m_{\rm proton}$) & $\delta_{\rm neo}'$ & $\delta_{\rm neo}'=\sum_s m_s n_s \frac{\partial \uparaDZ}{\partial\rho}$ \\
\tt SGNEO 	& kg/(m$^2$ s T)/($10^{19}$ $m_{\rm proton}$) & $\sigma_{\rm neo}$ &  $\sigma_{\rm neo}=\sum_s m_s n_s K_s$\\
\tt SGNEOD 	& kg/(m$^3$ s T)/($10^{19}$ $m_{\rm proton}$) & $\sigma_{\rm neo}'$ &  $\sigma_{\rm neo}'=\sum_s m_s n_s \frac{\partial K_s}{\partial\rho}$\\
\tt DDNEO 	& kg/(m$^2$ s)/($10^{19}$ $m_{\rm proton}$) & $\delta_{\rm d}$ & $\delta_{\rm d}= \sum_s m_s n_s  I H_{\rm p,s}/B_0$ \\
\tt UPAR 	& m/s	& $\upara$ & Parallel velocity as defined in equation (\ref{uparfsadef}) \\
\tt XUPAR 	& m$^2$/s 	&$\chi$ & Radial diffusivity of parallel momentum \\
\tt XUPAP 	& m$^2$/s 	&$\chi_\theta$ & Radial diffusivity of poloidal momentum \\
\tt CNPAR 	& m/s  	& $V_$  & Radial pinch of parallel momentum \\
\tt CNPAD 	& m/s  	& $V_d$  & Radial pinch of diamagnetic momentum \\
\tt CNPAP 	& m/s  	& $V_\theta$  & Radial pinch of poloidal momentum \\
\tt RUPAR 	& 1/s$^2$ 	&  $R_{\parallel}$ &  Parallel Reynold stress $R_{\parallel}=\frac{\sum_s\left\langle R \frac{\bphi}{B} \Pi_{\parallel\Psi,s} 
|\nabla{\rho}| \right\rangle }{ (10^{19} m_{\rm proton})}$\\
\tt RUPYR 	& 1/s$^2$  	& $R_y$ & Perp Reynold stress $R_y=\frac{\sum_s\left\langle R \frac{\bpol}{B} \Pi_{y\Psi,s}
|\nabla{\rho}| \right\rangle}{(10^{19} m_{\rm proton})}$\\
\tt RUPFR 	& 1/s$^2$  	& $R_{\phi}$ & Toroidal Reynold stress $R_{\phi}=\frac{\sum_s\left\langle R  \Pi_{\phi\Psi,s}
|\nabla{\rho}| \right\rangle}{(10^{19} m_{\rm proton})}$\\
\tt TTRQ 	& N/m$^2$ /($10^{19}$ $m_{\rm proton}$)   	& $\left\langle R \sum_s S_{\phi,s}\right\rangle$ & Toroidal input torque \\
\hline
\end{tabular}
\bigskip\\
Notice that MRHO, DLNEO, SGNEO, XUPAR, XUPAP, CNPAR, CNPAD, CNPAP, RUPAR, RUPYR, RUPFR, TTRQ, DLNEOD, DDNEO, SGNEOD have to be provided by the user in the model file, while UPAR is the evolved quantity.
%
\subsection{Dedicated control parameters}
\\
{\bf Table 2.2} Control parameter {\bf IPROT}\\[2ex]
\begin{tabular}{|l|l|l|}					\hline
\bf Name	&\bf Value & Description\\ \hline
\tt IPROT 	& 0 & Include poloidal rotation terms \\
\tt  	& 1 &  Neglect poloidal rotation terms \\ 
\hline
\end{tabular}
\bigskip\\
By default, {\bf IPROT = 0}.


\clearpage
